\chapter{Optional Stopping and Random Walks}
\label{sec:p15-optional-stopping-and-random-walks}

Stopping at a bounded time preserves a martingale's mean.
At an unbounded time, the conclusion requires control of the
stopped values or of the increments accumulated before stopping.
This chapter supplies the complete optional-stopping presentation
and its random-walk exercises in Durrett's Section 4.8
\cite[pp.~255--262]{DurrettPTEAuthor2019}.

\begin{hbdefinition}{def:p15-os-conventions}{Stopped values, infinite times, and terminal limits}
Work on \((\Omega,\mathcal F,\mathbb P)\) with a filtration
\((\mathcal F_n)_{n\geq0}\). Processes are real and adapted.
A submartingale, supermartingale or martingale is integrable at
every finite time, with its respective conditional inequality or
equality: \(\mathbb E(X_{n+1}\mid\mathcal F_n)\geq X_n\),
\(\leq X_n\), or \(=X_n\), respectively.
All conditional identities and inequalities are almost-sure statements.
A stopping time takes values in
\(\mathbb N_0\cup\{\infty\}\), where \(\mathbb N_0=\{0,1,\ldots\}\);
it satisfies \(\{T\leq n\}\in\mathcal F_n\) for every \(n\).
The convention is \(\inf\varnothing=\infty\). Its stopped information is
\[
 \mathcal F_T=\{A\in\mathcal F:
                    A\cap\{T\leq n\}\in\mathcal F_n
                    \text{ for every }n\}.
 \]
Whenever a finite terminal limit is specified on \(\{T=\infty\}\),
define \(X_T=X_j\) on \(\{T=j<\infty\}\) and use that limit
on \(\{T=\infty\}\). Limits are assigned zero on their measurable
null exceptional sets. If \(T\) is almost surely finite, this
simply assigns zero on \(\{T=\infty\}\).
Every use of \(X_T\) at a genuinely infinite time states the
required limiting hypothesis.

A family of integrable variables \(Z_i\) is uniformly integrable
when
\(\sup_i\mathbb E(|Z_i|;|Z_i|>K)\to0\) as \(K\to\infty\).
Write \(X_n^+=\max(X_n,0)\), \(X_n^-=\max(-X_n,0)\), and
\(\mathbb E(Z;A)=\mathbb E(Z\mathbf1_A)\).
\end{hbdefinition}

\section{Sampling with an integrable terminal value}
\label{sec:p15-os-sampling}

\begin{hbtheorem}{thm:p15-os-ui-stopped}{proved=annex:p15.os-ui-stopped}{Stopping a uniformly integrable process}
If \(X_n\) is a uniformly integrable submartingale and
\(T\leq\infty\) is any stopping time, then
\((X_{T\wedge n})_{n\geq0}\) is uniformly integrable and
\[
 X_{T\wedge n}\to X_T\quad\text{almost surely and in }L^1.
 \]
Here \(X_\infty=\lim_nX_n\) exists, is finite and integrable,
and supplies \(X_T\) on \(\{T=\infty\}\).
Moreover
\[
 \mathbb EX_0\leq\mathbb EX_T\leq\mathbb EX_\infty.
 \]
For supermartingales the inequalities reverse; for martingales
they are equalities. The annex retains the stopped-positive-part
and tail-splitting proof of Theorems 4.8.1 and 4.8.3.
\end{hbtheorem}

\begin{hbtheorem}{thm:p15-os-residual-ui}{proved=annex:p15.os-residual-ui}{The residual-tail criterion}
Let \(X_n\) be a submartingale and \(T\) a stopping time.
Assume \(X_T\) is defined by the preceding convention,
\(\mathbb E|X_T|<\infty\), and
\((X_n\mathbf1_{\{T>n\}})_{n\geq0}\) is uniformly integrable.
When \(\mathbb P(T=\infty)>0\), the convention here requires
\(X_n\to X_T\) almost surely on that event.
Then \(X_{T\wedge n}\) is uniformly integrable,
\(X_{T\wedge n}\to X_T\) in \(L^1\), and
\(\mathbb EX_0\leq\mathbb EX_T\).
The sign-reversed and equality versions hold as well.
This is the complete limiting premise of source Theorem 4.8.2.
\end{hbtheorem}

\begin{hbtheorem}{thm:p15-os-conditional-sampling}{proved=annex:p15.os-conditional-sampling}{Conditional sampling at two unbounded times}
Let \(L\leq M\leq\infty\) be stopping times and \(Y\) a
submartingale. If the stopped process \(Y_{M\wedge n}\) is
uniformly integrable, then its almost-sure and \(L^1\) limit
defines \(Y_M\), and \(Y_L\) is integrable and
\(\mathcal F_L\)-measurable. One has
\[
 Y_L\leq\mathbb E(Y_M\mid\mathcal F_L),\qquad
 \mathbb EY_L\leq\mathbb EY_M.
 \]
On \(\{L=\infty\}\), both times are infinite and both sampled
values use the same terminal limit.
For supermartingales reverse the inequalities; for martingales
use equalities. This supplies Exercise 4.8.1 with genuinely
unbounded and possibly infinite stopping times.
\end{hbtheorem}

\begin{hbtheorem}{thm:p15-os-nonnegative}{proved=annex:p15.os-nonnegative}{Nonnegative supermartingales and their crossing bound}
For a nonnegative supermartingale \(X\), its almost-sure limit
\(X_\infty\) is finite and integrable. For stopping times
\(L\leq M\leq\infty\),
\[
 \mathbb E(X_M\mid\mathcal F_L)\leq X_L,\qquad
 \mathbb EX_M\leq\mathbb EX_L\leq\mathbb EX_0.
 \]
In particular, for \(\lambda>0\),
\[
 \mathbb P\left(\sup_{n\geq0}X_n\geq\lambda\right)
                         \leq\frac{\mathbb EX_0}{\lambda};
 \]
the same bound holds with a strict crossing.
This includes Theorem 4.8.4 and Exercise 4.8.2.
\end{hbtheorem}

\begin{hbtheorem}{thm:p15-os-integrable-time}{proved=annex:p15.os-integrable-time}{Integrable time and a conditional increment bound}
Suppose \(X\) is a submartingale and a deterministic constant
\(0\leq B<\infty\) satisfies
\[
 \mathbb E(|X_{n+1}-X_n|\mid\mathcal F_n)\leq B
                                    \quad(n\geq0).
 \]
If \(T\) is a stopping time with \(\mathbb ET<\infty\), then
\(X_{T\wedge n}\) is dominated in absolute value by one
integrable variable, hence uniformly integrable, and
\[
 \mathbb E|X_T|\leq\mathbb E|X_0|+B\mathbb ET,\qquad
 \mathbb EX_T\geq\mathbb EX_0.
 \]
Again reverse signs for a supermartingale and use equality for
a martingale.
\end{hbtheorem}

\section{Wald identities and restarting the increments}
\label{sec:p15-os-wald}

\begin{hbdefinition}{def:p15-os-walk}{Random walks and probability generating functions}
Let \(\xi_1,\xi_2,\ldots\) be independent increments, and let
\[
 S_0=x,\quad S_n=x+\sum_{j=1}^n\xi_j,\qquad
 \mathcal F_n=\sigma(\xi_1,\ldots,\xi_n),
 \]
with trivial \(\mathcal F_0\) and deterministic real \(x\).
Identical laws will be stated where required.
\(\mathbb P_x,\mathbb E_x\) indicate the starting point \(x\).
In an i.i.d. model, write
\(\varphi(\theta)=\mathbb Ee^{\theta\xi_1}\) for a real
\(\theta\) whenever this exponential moment is finite.
This is an exponential moment, rather than a characteristic
function with an imaginary exponent.
For integer-valued increments, define
\(T_a=\inf\{n\geq0:S_n=a\}\).
If \(T\) is a nonnegative integer-valued, possibly infinite
time, its probability generating function for \(0\leq s<1\) is
\[
 G_T(s)=\mathbb E(s^T;T<\infty),
 \]
with \(0^0=1\). At \(s=1\) set \(G_T(1)=\mathbb P(T<\infty)\).
When \(T<\infty\) almost surely this is the usual generating
function. The convention never defines \(1^\infty\) as an
unjustified terminal value.
\end{hbdefinition}

\begin{hbtheorem}{thm:p15-os-wald-first}{proved=annex:p15.os-wald-first}{Wald's first identity}
If the increments are i.i.d., \(\mathbb E|\xi_1|<\infty\),
\(\mu=\mathbb E\xi_1\), and \(T\) is a stopping time with
\(\mathbb ET<\infty\), then \(S_T\) is integrable and
\[
 \mathbb E_xS_T=x+\mu\mathbb E_xT.
 \]
This includes Theorem 4.8.6 at \(x=0\).
\end{hbtheorem}

\begin{hbtheorem}{thm:p15-os-wald-second}{proved=annex:p15.os-wald-second}{Wald's second identity and an exit-time lower bound}
Let the increments be independent, with
\(\mathbb E\xi_j=0\) and
\(\mathbb E\xi_j^2=\sigma^2\in[0,\infty)\) for every \(j\);
identical laws are unnecessary. With \(S_0=0\), if
\(\mathbb ET<\infty\), then
\[
 S_{T\wedge n}\to S_T\text{ in }L^2,\qquad
 \mathbb ES_T^2=\sigma^2\mathbb ET.
 \]
The proof establishes the stopped-square integrability before
passing to the limit in the quadratic martingale.
For \(\sigma^2>0\), \(a>0\), and
\(T=\inf\{n\geq0:|S_n|>a\}\), one has
\[
 \mathbb ET\geq a^2/\sigma^2
 \]
as an extended nonnegative inequality. If \(\mathbb ET<\infty\),
the inequality is strict. These are Exercises 4.8.3--4.8.4.
When \(\sigma^2=0\), the centered increments are all zero
almost surely; the exit time from \([-a,a]\), \(a>0\), is
infinite, and the ratio with \(\sigma^2\) is not used.
\end{hbtheorem}

\begin{hblemma}{thm:p15-os-independent-restart}{proved=annex:p15.os-independent-restart}{Independent increments after a finite stopping time}
For i.i.d. increments and an almost-surely finite stopping time \(T\),
the sequence \(\xi_{T+1},\xi_{T+2},\ldots\) has the original
i.i.d. law and is independent of \(\mathcal F_T\).
Values on \(\{T=\infty\}\) may be assigned from any fixed
increment sequence, since that event is null.
In particular the successive one-level downward passage times
of a nearest-neighbor walk with negative drift are independent
and have the same law.
\end{hblemma}

\section{Nearest-neighbor exits and passage times}
\label{sec:p15-os-simple-walks}

\begin{hbdefinition}{def:p15-os-simple-walk}{The lattice and the two directions of drift}
A simple walk has i.i.d. increments with
\(\mathbb P(\xi_j=1)=p\), \(\mathbb P(\xi_j=-1)=q=1-p\).
For the nondeterministic model \(0<p<1\).
It is symmetric when \(p=q=1/2\), and biased otherwise.
For integers \(a<x<b\), let
\(\tau_{a,b}=\inf\{n\geq0:S_n\notin(a,b)\}=T_a\wedge T_b\).
The integer endpoints ensure that the exit value is exactly \(a\)
or \(b\); a general walk can overshoot.
When \(p\ne q\), write \(r=q/p>0\) and \(d=p-q\ne0\).
\end{hbdefinition}

\begin{hblemma}{thm:p15-os-finite-exit}{proved=annex:p15.os-finite-exit}{An exit bound that establishes all time moments first}
For \(0<p<1\) and the preceding interval, put \(w=b-a\).
For every integer \(m\geq0\),
\[
 \mathbb P_x(\tau_{a,b}>mw)\leq(1-p^w)^m.
 \]
Thus the exit is almost surely finite, has every positive
integer moment finite, and has a finite exponential moment
for every \(c>0\) with \(e^{cw}(1-p^w)<1\).
At \(p=1\) the exit time is \(b-x\) and the upper boundary
is certain; at \(p=0\) it is \(x-a\) and the lower boundary
is certain. A starting point at either boundary has exit time zero.
\end{hblemma}

\begin{hbtheorem}{thm:p15-os-fair-exit}{proved=annex:p15.os-fair-exit}{Fair gambler's ruin: probabilities and mean time}
For the symmetric walk and integers \(a<x<b\),
\[
 \mathbb P_x(S_{\tau_{a,b}}=a)=\frac{b-x}{b-a},\qquad
 \mathbb P_x(S_{\tau_{a,b}}=b)=\frac{x-a}{b-a},\qquad
 \mathbb E_x\tau_{a,b}=(b-x)(x-a).
 \]
In particular \(\mathbb E_0\tau_{a,b}=-ab\) when \(a<0<b\).
Every integer level is hit almost surely.
The proof uses the linear martingale and the bounded-time
quadratic martingale \(S_n^2-n\), retaining the source limits.
\end{hbtheorem}

\begin{hbtheorem}{thm:p15-os-fair-hitting}{proved=annex:p15.os-fair-hitting}{The one-level generating function and its full inversion}
For the symmetric walk with \(S_0=0\), let \(T_1\) be its
first hit of \(1\). Then
\[
 \mathbb E_0s^{T_1}=\frac{1-\sqrt{1-s^2}}s\quad(0<s<1),
 \]
with endpoint values \(0\) at \(s=0\) and \(1\) at \(s=1\).
For integers \(n\geq1\),
\[
 \mathbb P_0(T_1=2n-1)
   =\frac1{2n-1}\binom{2n}{n}2^{-2n},
 \qquad \mathbb P_0(T_1\text{ is even})=0.
 \]
Moreover \(\mathbb E_0T_1=\infty\).
The source exponential-martingale and binomial-series proof
is supplied with the distinct roots for \(e^\theta\) and
its reciprocal kept explicit.
\end{hbtheorem}

\begin{hbtheorem}{thm:p15-os-biased-exit}{proved=annex:p15.os-biased-exit}{Biased gambler's ruin}
For \(0<p<1\), \(p\ne q\), the process \(r^{S_n}\) is a
martingale. For integers \(a<x<b\),
\[
 \alpha=\mathbb P_x(T_a<T_b)=\frac{r^b-r^x}{r^b-r^a},
 \qquad
 \beta=\mathbb P_x(T_b<T_a)=\frac{r^x-r^a}{r^b-r^a},
 \qquad
 \mathbb E_x\tau_{a,b}=\frac{a\alpha+b\beta-x}{p-q}.
 \]
For the source roulette model \(p=18/38\), starting at \(20\),
the probability of reaching \(40\) before zero is
\[
 \frac{(10/9)^{20}-1}{(10/9)^{40}-1}
       =\frac1{(10/9)^{20}+1}\approx0.1084.
 \]
This includes Theorem 4.8.9(a)--(b) and Example 4.8.10.
\end{hbtheorem}

\begin{hbtheorem}{thm:p15-os-drift-hitting}{proved=annex:p15.os-drift-hitting}{One-sided hitting and its mean without a circular moment assumption}
For \(1/2<p<1\), \(S_0=0\), and integers \(a<0<b\),
\[
 \mathbb P_0(T_a<\infty)=r^{-a},\qquad
 \mathbb P_0(T_b<\infty)=1,\qquad
 \mathbb E_0T_b=\frac b{p-q}.
 \]
The depth \(I=-\inf_{n\geq0}S_n\) has
\(\mathbb P(I\geq k)=r^k\), \(k\geq1\), and
\(\mathbb EI=r/(1-r)<\infty\).
This supplies the integrable lower bound needed to prove the
mean hitting time.
By reflection, for \(0<p<1/2\) and integer \(x\geq1\),
\[
 V_0=T_0<\infty\ \text{a.s.},\qquad
\mathbb E_xV_0=x/(q-p).
\]
At \(p=1\), the walk increases deterministically:
\(T_b=b\) from zero and a lower level is never hit.
At \(p=0\), the reflected downward assertions are deterministic.
\end{hbtheorem}

\begin{hbtheorem}{thm:p15-os-downward-variance}{proved=annex:p15.os-downward-variance}{Variance of the time to ruin}
For \(0<p<1/2\), \(x\geq1\) integer, and
\(\rho=q-p>0\),
\[
 \operatorname{Var}_x(V_0)
       =\frac{x(1-\rho^2)}{\rho^3}
       =\frac{4pq\,x}{(q-p)^3}.
 \]
The second moment is finite and equals this variance plus
\(x^2/\rho^2\). Finiteness follows during the Wald-second
calculation, before the time variance is expanded.
The law of \(V_0\) is a sum of \(x\) independent identically
distributed one-level passage times, explaining the factor \(x\).
At \(p=0\), \(V_0=x\) deterministically, so its variance is zero,
as the displayed expression also gives. At \(x=0\), the time,
mean, and variance are all zero.
\end{hbtheorem}

\begin{hbtheorem}{thm:p15-os-downward-pgf}{proved=annex:p15.os-downward-pgf}{Biased ruin generating functions and coefficients}
In the same negative-drift model,
\(\varphi(\theta)=pe^\theta+qe^{-\theta}\).
For \(\theta\leq0\),
\[
 e^{\theta x}=\mathbb E_x\{\varphi(\theta)^{-V_0}\}.
 \]
For \(0<s<1\),
\[
 \mathbb E_xs^{V_0}=f(s)^x,\qquad
 f(s)=\frac{1-\sqrt{1-4pqs^2}}{2ps}.
 \]
For \(x=1\) the complete inversion is
\[
 \mathbb P_1(V_0=2k+1)=C_kp^kq^{k+1},\qquad
 C_k=\frac1{k+1}\binom{2k}{k},\quad k\geq0.
 \]
For general \(x\geq1\), the mass at \(x+2k\) is
\[
 p^kq^{k+x}
     \sum_{\substack{k_1+\cdots+k_x=k\\k_j\geq0}}
                                 \prod_{j=1}^x C_{k_j};
 \]
all other times have zero mass.
The endpoints are \(G_{V_0}(0)=0\), \(G_{V_0}(1)=1\).
At \(x=0\) the generating function is identically one.
At \(p=0\), \(V_0=x\) deterministically and its function is \(s^x\).
\end{hbtheorem}

\begin{hbtheorem}{thm:p15-os-quartic-time}{proved=annex:p15.os-quartic-time}{The quartic martingale and a two-sided exit-time second moment}
For the symmetric walk from zero,
\[
 Y_n=S_n^4-6nS_n^2+3n^2+2n
 \]
is a martingale. For integer \(a\geq1\) and
\(T=\inf\{n\geq0:S_n\notin(-a,a)\}\),
\[
 \mathbb ET=a^2,\qquad
 \mathbb ET^2=\frac{5a^4-2a^2}{3},\qquad
 \operatorname{Var}(T)=\frac{2a^2(a^2-1)}3.
 \]
The proof establishes the necessary time integrability before
passing to the unbounded stopped identity.
\end{hbtheorem}

\section{Exponential bounds and exact downward ruin}
\label{sec:p15-os-exponential-ruin}

\begin{hbtheorem}{thm:p15-os-exponential-exit}{proved=annex:p15.os-exponential-exit}{An exponential exit identity with unbounded overshoots}
Let the increments be i.i.d., nonconstant, and suppose
\(\theta_0<0\) satisfies
\(\varphi(\theta_0)=\mathbb Ee^{\theta_0\xi_1}=1\).
For real \(a<x<b\), define
\(\tau=\inf\{n\geq0:S_n\notin(a,b)\}\).
Then \(\mathbb E_x\tau<\infty\),
\(\mathbb E_x e^{\theta_0S_\tau}<\infty\), and
\[
 \mathbb E_x e^{\theta_0S_\tau}=e^{\theta_0x},\qquad
 \mathbb P_x(S_\tau\leq a)\leq e^{\theta_0(x-a)}.
 \]
No bounded overshoot assumption is imposed.
For the possibly infinite time
\(T=\inf\{n\geq0:S_n\leq a\}\), one also has
\[
 \mathbb P_x(T<\infty)\leq e^{\theta_0(x-a)}.
 \]
At \(x=0\) these are the source Exercise 4.8.8 formulas.
\end{hbtheorem}

\begin{hbtheorem}{thm:p15-os-skip-free-ruin}{proved=annex:p15.os-skip-free-ruin}{Exact ruin when a downward step cannot skip a level}
Suppose the increments are i.i.d. integers with
\(\xi_1\geq-1\), \(\mathbb P(\xi_1=-1)>0\), and
\(\mathbb E\xi_1\in(0,\infty]\).
There is a unique \(\theta_0<0\) with
\(\mathbb Ee^{\theta_0\xi_1}=1\).
For integers \(a<x\),
\[
 \mathbb P_x(T_a<\infty)=e^{\theta_0(x-a)}.
 \]
The proof handles the infinite stopping event explicitly;
it uses the positive-drift strong law also when the positive
mean is infinite.
For the favorable game with increments \(-1,1,2\), each of
probability \(1/3\), the ruin probability from integer
capital \(i\geq1\) is
\[
                    (\sqrt2-1)^i.
 \]
At initial capital zero it is one.
\end{hbtheorem}

\begin{hbtheorem}{thm:p15-os-insurance}{proved=annex:p15.os-insurance}{The Gaussian-claim ruin bound}
Let independent claims \(\zeta_n\) have the common normal law
\(N(\mu,\sigma^2)\), with \(c>\mu\), \(c>0\), and
\(\sigma^2>0\). For deterministic \(S_0=x>0\) let
\(S_n=x+\sum_{j=1}^n(c-\zeta_j)\).
Ruin means \(S_n\leq0\) at some time. Then
\[
 \mathbb P_x(\text{ruin})
              \leq\exp\{-2(c-\mu)x/\sigma^2\}.
 \]
If \(\sigma^2=0\), the assets increase deterministically and
the ruin probability is zero; the displayed ratio is not
used in that case.
\end{hbtheorem}

\begin{hbtheorem}{thm:p15-os-boundaries}{proved=annex:p15.os-boundaries}{Which stopping hypotheses cannot be omitted}
The following boundary examples have complete proofs.
\begin{enumerate}
\item For the fair walk from zero, \(T_1<\infty\) almost surely
 and \(S_{T_1}=1\), but \(\mathbb ES_0=0\) and
 \(\mathbb ET_1=\infty\). Integrability of the terminal value
 alone does not justify equality of means.
\item For the fair walk from one stopped at zero,
 the nonnegative martingale has mean one at every finite horizon
 and terminal mean zero. The nonnegative supermartingale inequality
 can be strict even when the original process is a martingale.
\item Doubling a fair stake until the first success gives an
 integrable stopping time of mean two and terminal gain one,
 while every finite stopped-horizon mean is zero. The conditional
 increment bound in Theorem~\ref{thm:p15-os-integrable-time}
 fails.
\item A bounded random time chosen using future information need
 not preserve a martingale's mean. With a fair sign \(\xi\),
 \(M_0=0,M_n=\xi\) for \(n\geq1\), and \(T=1\) on \(\{\xi=1\}\), \(T=0\)
 otherwise, \(\mathbb EM_T=1/2\); this \(T\) is not a stopping time
 for the trivial initial information.
\end{enumerate}
\end{hbtheorem}

\begin{proofboundary}{def:p15-os-next-boundary}{The next distinct source presentation}
This chapter supplies Theorems 4.8.1--4.8.9, Example 4.8.10,
and Exercises 4.8.1--4.8.11, including the possibly infinite
sampling times where stated.
Section 4.9's reflection, ballot, return-count and arcsine arguments,
together with both parts of its renewal exercise, are supplied in
\cref{sec:p15-reflection-and-ballot-counts}. The path-counting proofs
there give a separate derivation of the return and passage laws.
The exact relation between successive returns and positive level hits
is \cref{thm:p15-rb-return-scaling}.
\end{proofboundary}
